\begin{thebibliography}{99}\small

\bibitem[Baitairian et~al.(2025)]{baitairian2025} Baitairian, J.-B., Sebastien, B., Jreich, R., Katsahian, S. and Guilloux, A. (2025). Calibrating confounding strength in sensitivity models for weighting estimators: a comparative review and a new method. arXiv:2510.16560.

\bibitem[Bickel et~al.(2009)]{bickel2009} Bickel, S., Brückner, M. and Scheffer, T. (2009). Discriminative learning under covariate shift. \emph{Journal of Machine Learning Research}, 10, 2137--2155.

\bibitem[Blake et~al.(2020)]{blake2020} Blake, H. A., Leyrat, C., Mansfield, K. E., Seaman, S., Tomlinson, L. A., Carpenter, J. and Williamson, E. J. (2020). Propensity scores using missingness pattern information: a practical guide. \emph{Statistics in Medicine}, 39, 1641--1657.

\bibitem[Chernozhukov et~al.(2018)]{chernozhukov2018} Chernozhukov, V., Chetverikov, D., Demirer, M., Duflo, E., Hansen, C., Newey, W. and Robins, J. (2018). Double/debiased machine learning for treatment and structural parameters. \emph{The Econometrics Journal}, 21(1), C1--C68.

\bibitem[Cook et~al.(2008)]{cook2008} Cook, T. D., Shadish, W. R. and Wong, V. C. (2008). Three conditions under which experiments and observational studies produce comparable causal estimates. \emph{Journal of Policy Analysis and Management}, 27(4), 724--750.

\bibitem[Crump et~al.(2009)]{crump2009} Crump, R. K., Hotz, V. J., Imbens, G. W. and Mitnik, O. A. (2009). Dealing with limited overlap in estimation of average treatment effects. \emph{Biometrika}, 96(1), 187--199.

\bibitem[D'Agostino and Rubin(2000)]{dagostino2000} D'Agostino, R. B., Jr. and Rubin, D. B. (2000). Estimating and using propensity scores with partially missing data. \emph{Journal of the American Statistical Association}, 95, 749--759.

\bibitem[D'Amour et~al.(2021)]{damour2021} D'Amour, A., Ding, P., Feller, A., Lei, L. and Sekhon, J. (2021). Overlap in observational studies with high-dimensional covariates. \emph{Journal of Econometrics}, 221, 644--654.

\bibitem[de~Luna et~al.(2011)]{deluna2011} de Luna, X., Waernbaum, I. and Richardson, T. S. (2011). Covariate selection for the nonparametric estimation of an average treatment effect. \emph{Biometrika}, 98(4), 861--875.

\bibitem[Dong et~al.(2020)]{dong2020} Dong, J., Zhang, J. L., Zeng, S. and Li, F. (2020). Subgroup balancing propensity score. \emph{Statistical Methods in Medical Research}, 29(3), 659--676.

\bibitem[Dorn and Guo(2023)]{dorn2023} Dorn, J. and Guo, K. (2023). Sharp sensitivity analysis for inverse propensity weighting via quantile balancing. \emph{Journal of the American Statistical Association}, 118(544), 2645--2657.

\bibitem[Hahn(1998)]{hahn1998} Hahn, J. (1998). On the role of the propensity score in efficient semiparametric estimation of average treatment effects. \emph{Econometrica}, 66(2), 315--332.

\bibitem[Heckman et~al.(1997)]{heckman1997} Heckman, J. J., Ichimura, H. and Todd, P. E. (1997). Matching as an econometric evaluation estimator: evidence from evaluating a job training programme. \emph{The Review of Economic Studies}, 64(4), 605--654.

\bibitem[Hirano et~al.(2003)]{hirano2003} Hirano, K., Imbens, G. W. and Ridder, G. (2003). Efficient estimation of average treatment effects using the estimated propensity score. \emph{Econometrica}, 71(4), 1161--1189.

\bibitem[Khan and Tamer(2010)]{khan2010} Khan, S. and Tamer, E. (2010). Irregular identification, support conditions, and inverse weight estimation. \emph{Econometrica}, 78(6), 2021--2042.

\bibitem[Liu et~al.(2023)]{liu2023} Liu, Y., Li, H., Zhou, Y. and Matsouaka, R. A. (2023). Average treatment effect on the treated, under lack of positivity. arXiv:2309.01334.

\bibitem[Ma and Namkoong(2025)]{manamkoong2025} Ma, Y. and Namkoong, H. (2025). A sensitivity approach to causal inference under limited overlap. arXiv:2511.22003.

\bibitem[Ma and Wang(2020)]{ma2020} Ma, X. and Wang, J. (2020). Robust inference using inverse probability weighting. \emph{Journal of the American Statistical Association}, 115, 1851--1860.

\bibitem[Ma et~al.(2023)]{ma2023} Ma, Y., Sant'Anna, P. H. C., Sasaki, Y. and Ura, T. (2023). Doubly robust estimators with weak overlap. arXiv:2304.08974.

\bibitem[McClean et~al.(2024)]{mcclean2024} McClean, A., Branson, Z. and Kennedy, E. H. (2024). Calibrated sensitivity models. arXiv:2405.08738.

\bibitem[Petersen et~al.(2012)]{petersen2012} Petersen, M. L., Porter, K. E., Gruber, S., Wang, Y. and van der Laan, M. J. (2012). Diagnosing and responding to violations in the positivity assumption. \emph{Statistical Methods in Medical Research}, 21(1), 31--54.

\bibitem[Rassen et~al.(2012)]{rassen2012} Rassen, J. A., Glynn, R. J., Rothman, K., Setoguchi, S. and Schneeweiss, S. (2012). Applying propensity scores estimated in a full cohort to adjust for confounding in subgroup analyses. \emph{Pharmacoepidemiology and Drug Safety}, 21(7), 697--709.

\bibitem[Rudin(1987)]{rudin1987} Rudin, W. (1987). \emph{Real and Complex Analysis} (3rd ed.). McGraw-Hill.

\bibitem[Steiner et~al.(2011)]{steiner2011} Steiner, P. M., Cook, T. D. and Shadish, W. R. (2011). On the importance of reliable covariate measurement in selection bias adjustments using propensity scores. \emph{Journal of Educational and Behavioral Statistics}, 36(2), 213--236.

\bibitem[Sugiyama et~al.(2012)]{sugiyama2012} Sugiyama, M., Suzuki, T. and Kanamori, T. (2012). \emph{Density Ratio Estimation in Machine Learning}. Cambridge University Press.

\bibitem[Tan(2006)]{tan2006} Tan, Z. (2006). A distributional approach for causal inference using propensity scores. \emph{Journal of the American Statistical Association}, 101(476), 1619--1634.

\bibitem[VanderWeele and Shpitser(2011)]{vanderweele2011} VanderWeele, T. J. and Shpitser, I. (2011). A new criterion for confounder selection. \emph{Biometrics}, 67(4), 1406--1413.

\bibitem[Westreich et~al.(2011)]{westreich2011} Westreich, D., Cole, S. R., Funk, M. J., Brookhart, M. A. and Stürmer, T. (2011). The role of the c-statistic in variable selection for propensity score models. \emph{Pharmacoepidemiology and Drug Safety}, 20(3), 317--320.

\end{thebibliography}
